\section{Information-Theoretic Consistency Check}

This section is not a proof of the empirical mechanism. Its role is narrower:
it checks that the observed non-identifiability is compatible with a simple
information view of retrieval.

Let $Y$ denote the answer-supporting information needed for query $q$, and let
each passage $p_i$ contain two kinds of signal: local query evidence measured
by a score $s(q,p_i)$, and answer-supporting information $I(Y;p_i\mid q)$ that
may not be captured by the local score. A refinement rule that keeps passages
or sub-passages only through $s(q,p_i)$ is using an incomplete statistic unless
$s$ is sufficient for answer support.

\begin{proposition}[Incomplete local scores]
If two passages can have the same local score but different
$I(Y;p_i\mid q)$, then a policy that ranks only by local score can replace a
more answer-informative context with a less answer-informative one while
leaving the reported retrieval score unchanged or improved.
\end{proposition}

\paragraph{Implication.}
The proposition is structural. It says that local scoring leaves room for
discarding answer-relevant information; it does not say that this failure is
common, large, or caused by CCS. Those are empirical questions, and the matched
coherence experiment rejects the strong causal version for the primary
SQuAD/Mistral setting.

\begin{theorem}[Consistency of the controlled tests]
Suppose a RAG claim is based on a statistic $z(q,P)$ that is not sufficient for
answer support. Then there can exist two context-selection policies with the
same or better value of $z$ but different answer faithfulness. Identifying the
faithfulness effect therefore requires additional controls, such as matching
local similarity, varying the evaluator, checking threshold transfer, or
reporting deployment cost.
\end{theorem}

\paragraph{Use in this paper.}
The theorem motivates the experimental design, not the conclusion. The results
come from the preregistered controls in Section~3. The information-theoretic
view only explains why a local retrieval score can be an insufficient basis for
faithfulness claims.
